\section{Einfluss der Zahnfußgeometrie}
\label{sec:fussgeometrie}

Die drei Fußrundungsformen aus \secref{sec:werkzeugkette} wurden mit den Parametern ihrer Quellen und mit den Optima der ebenen Parameterstudie am Rad gerechnet. Jede Variante ersetzt die Trochoide der Erzeugung zwischen dem Übergangsdurchmesser und der Lückenmitte. Netzvorlage, Ritzelfläche, Lastschritte und Stellungen bleiben die der Taschen-Grundlinie aus \secref{sec:fe_modell}. Gerechnet wurden die Werkstoffstufen W1 bis W4 mit den Dehnratenkarten QS, DY1 und DY2 an den Normpunkten A bis E, je Variante 60 Rechnungen. Alle sechs Varianten liegen vollständig vor.

\tabref{tab:fussrundung_konturen} stellt die Konturen nebeneinander. Die Ellipse nach Frühe \cite{Frue12} beginnt am Fußformkreis und legt den Fußkreis um $0{,}57$\,mm tiefer als die Trochoide. Ihr kleinster Krümmungsradius liegt bei $0{,}083$\,mm in der Lückenmitte. Die beiden Patentkonturen setzen am mittleren Durchmesser zwischen Nutzkreis und Fußformkreis an und füllen den Fuß bis auf $50{,}2$\,mm und $50{,}1$\,mm auf. Dadurch rückt der kleinste Abstand der Ritzelkopfkante zur Fußrundung von $0{,}49$\,mm auf $0{,}20$\,mm. Der Vorauswahl diente eine ebene Parameterstudie am Einzahnsektor mit dem eigenen Löser. Sie trägt das ganze Moment auf einem Zahn mit einem Lastfleck am Normpunkt B und liefert ein relatives Maß (\secref{sec:werkzeugkette}).

% TK erledigt (eigener Löser in 3.6 beschrieben: ebener Dehnungszustand, CPE4/CPE4I-Formulierungen, Abgleich mit Abaqus): Mir ist gerade aufgefallen, das noch an keiner Stelle eine Aussage zum eigenen FE Löser gesagt wurde. Wie der aufgebaut ist oder welche Einschränkungen der hat. Das wäre ja wie ein eigenes Abaqus oder wie kann ich mir das vorstellen, darüber hatten wir auch noch nicht gesprochen und du hast den einfach implementiert auf welcher Basis aber?

\begin{table}[htbp]
	\caption{Fußrundungsvarianten am Rad kst-E: Kennwerte der Konturen und Änderung der größten Zugspannung gegen die Trochoide nach der ebenen Parameterstudie (W1, $16$\,Nm, Lastfleck an B).}
	\label{tab:fussrundung_konturen}
	\centering
	\footnotesize
	\setlength{\tabcolsep}{4pt}
	\begin{tabular}{L{3.0cm}L{2.9cm}lllll}
		Kontur                  & Parameter                                        & Übergang   & Fußkreis   & $\rho_{\mathrm{min}}$ & Abstand Kopfkante & eben      \\
		                        &                                                  & mm         & mm         & mm                    & mm                &           \\
		\hline
		\arrayrulecolor{black!30}
		Trochoide der Erzeugung & $\rho_{\mathrm{aP0}}$                            & $50{,}157$ & $49{,}524$ & $0{,}240$             & $0{,}490$         & $0$       \\
		\hline
		Ellipse nach Frühe      & $\gamma$ $30$°, $a/b$ $3{,}0$                    & $50{,}157$ & $48{,}951$ & $0{,}083$             & $0{,}491$         & $-13$\,\% \\
		\hline
		Ellipse, Optimum        & $\gamma$ $25$°, $a/b$ $2{,}0$                    & $50{,}157$ & $49{,}041$ & $0{,}170$             & $0{,}491$         & $-25$\,\% \\
		\hline
		Bézier nach Patent      & $f$ $1{,}0$                                      & $51{,}210$ & $50{,}215$ & $0{,}633$             & $0{,}203$         & $-32$\,\% \\
		\hline
		Bézier, Optimum         & $f$ $2{,}5$                                      & $51{,}210$ & $49{,}796$ & $0{,}366$             & $0{,}203$         & $-35$\,\% \\
		\hline
		Bionisch nach Patent    & $\varphi$ $24{,}1$°, $r$ $0{,}3\,S_{\mathrm{L}}$ & $51{,}210$ & $50{,}124$ & $0{,}382$             & $0{,}203$         & $-32$\,\% \\
		\hline
		Bionisch, Optimum       & $\varphi$ $35$°, $r$ $0{,}2\,S_{\mathrm{L}}$     & $51{,}210$ & $49{,}811$ & $0{,}255$             & $0{,}203$         & $-38$\,\% \\
		\hline
		\arrayrulecolor{black}
	\end{tabular}
\end{table}

% TK (folgt mit Paket E10: Eingriffsstellung C der drei eigenen Optima als Bilder pair_mesh in den Anhang): vielleicht kann man da auch nochmal jeweils Plots der Eingriffsstellung der unterschiedlichen Geometrien (die eigenen optimierten reichen) in den Anhang mit aufnhemn.

\subsection{Zugspannung im Zahnfuß}
\label{sec:fussgeometrie_sigma1}

\begin{table}[htbp]
	\caption{Änderung der größten Zugspannung $\sigma_{\mathrm{1}}$ und des Verdrehwinkels des Ritzels gegen die Trochoide am Normpunkt C, $16$\,Nm, quasistatische Karte, je Werkstoffstufe.}
	\label{tab:fussrundung_sigma1_c}
	\centering
	\begin{tabular}{llllll}
		Kontur               & W1            & W2            & W3            & W4            & Verdrehwinkel W1 \\
		\hline
		\arrayrulecolor{black!30}
		Ellipse nach Frühe   & $-15{,}9$\,\% & $-6{,}0$\,\%  & $-4{,}3$\,\%  & $-7{,}8$\,\%  & $+6{,}3$\,\%  \\
		\hline
		Ellipse, Optimum     & $-24{,}4$\,\% & $-13{,}4$\,\% & $-15{,}0$\,\% & $-12{,}3$\,\% & $+5{,}4$\,\%  \\
		\hline
		Bézier nach Patent   & $-25{,}2$\,\% & $-15{,}7$\,\% & $-25{,}0$\,\% & $-17{,}2$\,\% & $-11{,}0$\,\% \\
		\hline
		Bézier, Optimum      & $-27{,}6$\,\% & $-20{,}5$\,\% & $-29{,}1$\,\% & $-25{,}4$\,\% & $-5{,}4$\,\%  \\
		\hline
		Bionisch nach Patent & $-27{,}4$\,\% & $-15{,}7$\,\% & $-26{,}2$\,\% & $-19{,}0$\,\% & $-10{,}2$\,\% \\
		\hline
		Bionisch, Optimum    & $-25{,}4$\,\% & $-15{,}8$\,\% & $-22{,}6$\,\% & $-21{,}2$\,\% & $-6{,}4$\,\%  \\
		\hline
		\arrayrulecolor{black}
	\end{tabular}
\end{table}

\figref{fig:fusseffekt_sigma1_16nm_qs} zeigt die relative Änderung der größten Zugspannung $\sigma_{\mathrm{1}}$ der Fußrundung gegen die Trochoide an den fünf Normpunkten bei $16$\,Nm und quasistatischer Karte, je ein Feld für W1 bis W4. Alle Werte sind Größtwerte über die Zahnbreite. Zwischen den Punkten ist der Verlauf nicht gerechnet, deshalb stehen Marker ohne Linie. \tabref{tab:fussrundung_sigma1_c} nennt die Zahlen am Normpunkt C, an dem die Grundlinie ihre größte Fußspannung erreicht.

\begin{figure}[htbp]
	\centering
	\includegraphics[width=0.95\textwidth]{kap5/fusseffekt_sigma1_16nm_qs}
	\caption{Fußrundungsvarianten gegen die Trochoide der Erzeugung: relative Änderung der größten Zugspannung $\sigma_{\mathrm{1}}$ an den Normpunkten A bis E je Werkstoffstufe (W1 bis W4 von oben nach unten), Taschenkörper, $16$\,Nm, quasistatische Karte. Größtwerte über die Zahnbreite, Marker ohne Linie.}
	\label{fig:fusseffekt_sigma1_16nm_qs}
\end{figure}

% TK (ja, bleibt das Ziel; Paket E10: Zwillinge für 8 und 12 Nm der Bilder aus Kapitel 5 in den Anhang F): Und allgemein für alle Grafiken, nimmst du dann in den anhang auch jeweils die nicht eingebundenen Zwillinge anderer Drehmomentstufen oder ... auf, die in den Haupttext einfach zu überflutend wären, aber trotzdem natürlich relevante INformationen enthalten? Ist das noch das Ziel für den Anhang?

Im räumlichen Modell bleibt die Reihenfolge der ebenen Studie im Wesentlichen erhalten, die Beträge am Normpunkt C fallen kleiner aus. Unter W1 sinkt die größte Zugspannung bei der Ellipse nach Frühe um $16$\,\% und beim Optimum der Ellipse um $24$\,\%. Die Bézierkonturen erreichen $25$\,\% und $28$\,\%, die bionische Kontur nach den Patentwerten $27$\,\% und das bionische Optimum der ebenen Studie $25$\,\%. Am Normpunkt B, an dem das ebene Modell belastet war, treffen die beiden Patentkonturen nach den Literaturwerten und das Bézier-Optimum den Wert der ebenen Studie auf drei Prozentpunkte. Das bionische Optimum verfehlt ihn. Die ebene Studie sagte ihm mit $38$\,\% die größte Minderung voraus, im räumlichen Modell bleiben an B $24$\,\% und an C $25$\,\%, so dass es hinter die Bézierkontur mit $f = 2{,}5$ zurückfällt. Die ebene Studie mit einem Lastfleck an B überzeichnet den Nutzen des flachen Fußes, weil sie die Lastaufteilung und den Kontakt nicht kennt. Die Ellipsen weichen an B ab, weil ihr Größtwert nicht auf dem Flankenteil der Fußrundung liegt, sondern in der Lückenmitte.

An allen fünf Normpunkten senken die Patentkonturen die Zugspannung. Unter W1 liegen die Änderungen zwischen $-20$\,\% und $-35$\,\%, unter W3 zwischen $-17$\,\% und $-33$\,\%. Das bionische Optimum mit seinem kleinen Radius liegt am unteren Ende dieser Spannen. Die Druckspannung $\sigma_{\mathrm{3}}$ der gegenüberliegenden Fußrundung fällt um bis zu $44$\,\%. Der aufgefüllte Fuß macht den Zahn steifer. Der Verdrehwinkel des Ritzels nimmt um $4$\,\% bis $22$\,\% ab, die Kopfverschiebung um bis zu $26$\,\%. Dadurch drückt die Kante des Ritzelkopfes an den Normpunkten B und E stärker auf die Flanke des Rades. Der größte Kontaktdruck steigt dort bei den Patentwerten um $21$\,\% bis $30$\,\%, bei den Optima der ebenen Studie um bis zu $53$\,\% (Bézier $f = 2{,}5$ und bionisch $\varphi = 35$°). Die Flankenpressung unter der Ritzelkopfkante ist nach \secref{sec:waelzweg} schon bei der Trochoide die höchste Beanspruchung des Rades und nach \secref{sec:schadensbilder} der Ort der Einarbeitung im Versuch. Jede Fußrundung, die den Fuß auffüllt, entlastet den Fuß und belastet die Flanke, so dass der Nutzen am Fuß mit einer höheren Pressung an der Kante erkauft wird.

Anders verhalten sich die Ellipsen. Am Normpunkt C sinkt die Zugspannung. An den Doppeleingriffspunkten A, B, D und E steigt sie bei der Ellipse nach Frühe um bis zu $17$\,\% unter W1, $22$\,\% unter W2, $33$\,\% unter W3 und $20$\,\% unter W4. Die Druckspannung wächst dort um bis zu $30$\,\% unter W1 und W2 und um bis zu $51$\,\% unter W3. Beim Optimum der Ellipse bleiben die Änderungen an diesen Punkten unter W1 negativ. Unter W2 bis W4 erreichen sie bis zu $+6$\,\%, $+15$\,\% und $+7$\,\%. In beiden Fällen sitzt der Größtwert in der Lückenmitte auf dem kleinsten Krümmungsradius der Kontur, bei der Trochoide dagegen auf dem Flankenteil der Fußrundung. Der Radius von $0{,}083$\,mm liegt unter der Auflösung der Netzvorlage, so dass die Zahlen der Ellipsen an der Lückenmitte netzbegrenzt sind. Die Reihenfolge gegenüber den Patentkonturen bleibt davon unberührt. Mit dem tieferen Fuß wird der Zahn nachgiebiger, der Verdrehwinkel des Ritzels wächst um $5$\,\% bis $11$\,\%.

% TK erledigt (Satz zur Flankenpressung als Grenze im Absatz darüber, Aufnahme in 7.3): Kommt dann irgendwo in der Auswertung auch ein Satz dazu, dass eh schon die Flankenpressung sehr sehr hoch ist und limitierend und eine Optimierung des Fußes die Flanke noch mehr beansprucht. Das ist ja mehr oder weniger auch das Ergebnis dieser Fußoptimierungsfestellung, richtig? Ich schreibe es nur hier schonmal mit hin, dass ich es nicht vergesse, mit aufzunehmen.

\subsection{Werkstoffstufen, Dehnraten und Last}
\label{sec:fussgeometrie_stufen}

\figref{fig:fusseffekt_raten_16nm} ordnet die Änderung am Normpunkt C je Variante, Werkstoffstufe und Dehnratenkarte. Die Plastizität der isotropen Karte halbiert den Nutzen der Varianten. Unter W2 bleiben von den $25$\,\% bis $28$\,\% der Patentkonturen $16$\,\% bis $21$\,\%. Unter den anisotropen Stufen W3 und W4 bleibt ein Nutzen von $4$\,\% bis $29$\,\% erhalten, bei der Ellipse nach Frühe wächst er von W3 zu W4 sogar. Die Dehnrate verschiebt die Änderungen bei den Patentkonturen nach den Literaturwerten um höchstens zwei Prozentpunkte. Bei den Ellipsen und beim bionischen Optimum unter W3 nehmen die Minderungen von QS zu DY2 um vier bis sechs Prozentpunkte ab, weil die steifere Karte die Lastaufteilung zum Nachbarzahn verschiebt. Mit der Last nimmt der Nutzen der Patentkonturen nach den Literaturwerten ab. Die Bézierkontur mindert die Zugspannung unter W1 bei $8$\,Nm um $32$\,\% und bei $16$\,Nm um $25$\,\%, die bionische Kontur um $33$\,\% und $27$\,\%. Das bionische Optimum bleibt unter W1 bei $25$\,\% auf allen drei Stufen, unter W3 wächst sein Nutzen mit der Last von $17$\,\% auf $23$\,\%.

% TK erledigt (Bild gegen das Dehnratenbild fusseffekt_raten_16nm getauscht): Ich vermute, du hast eine falsche Grafik. Die Grafik beinhaltet gar nicht die Dehnraten sondern ist mehr oder weniger eine Dopplung von Bild 5.10 oder? nur eben über Normpunkte geplottet und dann allgemein. Das würde ich nicht machen bzw. eher dann eine Grafik mit den Dehnraten. Dann passt es auch zur aussage des Texts, da die Werkstoffstufen mehr oder weniger davor schon eingearbeitet werden. also bitte korrigieren.

\begin{figure}[htbp]
	\centering
	\includegraphics[width=0.95\textwidth]{kap5/fusseffekt_raten_16nm}
	\caption{Änderung der größten Zugspannung $\sigma_{\mathrm{1}}$ gegen die Trochoide am Normpunkt C je Variante, Werkstoffstufe (W1 bis W4 von oben nach unten) und Dehnratenkarte, Taschenkörper, $16$\,Nm. Marker ohne Linie.}
	\label{fig:fusseffekt_raten_16nm}
\end{figure}

Für den Nachweis nach Stommel et al. an der Fußrundung einer Variante ist die Lastgeschichte über alle Stellungen nötig. Die Grundserie des Taschenkörpers mit 99 Stellungen und W1 bis W4 bei quasistatischer Karte liegt vor und trägt den Nachweis in \secref{sec:stommel_ergebnis}. Die Vollserie des bionischen Optimums wurde zugunsten der Serien mit gemessener Geometrie zurückgestellt und bleibt Ausblick.

\subsection{Fertigungseinfluss der gemessenen Geometrie}
\label{sec:fussgeometrie_fertigung}

\begin{figure}[!htbp]
	\centering
	\includegraphics[width=0.95\textwidth]{kap5/kontur_zaehne_gemessen}
	\caption{Ritzel (links) und Rad (rechts) wie gefertigt gegen die STplus-Erzeugung, ganzer Zahn im Stirnschnitt in wahrem Maßstab. Mediane der GINA-Messung, Kopfrundung des Rades aus den Konturscans, Kopfrücknahme des Ritzels aus der GINA-Messung.}
	\label{fig:kontur_gemessen}
\end{figure}

Die Verzahnungsmessung in \secref{sec:pruefraeder_vorher} liefert die Abweichungen der gefertigten Teile von der Nenngeometrie, die Profil-Winkelabweichung, die Kopfrücknahme des Ritzels, die Kopfrundung des Rades von $0{,}21$\,mm und die Zahndicken. Daraus sind drei Rechenvarianten des Taschenkörpers aufgebaut: das Ritzel wie gemessen mit dem Rad der Nenngeometrie, das Rad wie gemessen mit dem Ritzel der Nenngeometrie und beide wie gemessen. Jede Variante trägt die Mediane der GINA-Messung über die Serie der Prüfstandsritzel und die Kernräder. Die Variante mit beiden Teilen wie gefertigt wurde als vollständige Serie über die 99 Stellungen mit W1 bis W4 und den drei Karten gerechnet. Das sind 1188 Rechnungen in $77{,}6$\,h mit drei Jobs zu je vier Kernen. Die beiden Einzelvarianten wurden an den Normpunkten A bis E gerechnet, je 60 Rechnungen in $3{,}7$ und $4{,}1$\,h. \figref{fig:kontur_gemessen} zeigt die gefertigten Zahnkonturen gegen die Erzeugung des STplus-Falls. Die Zahnweiten beider Teile treffen die STplus-Erzeugung der Nennserie, die das Zahndickenabmaß der Zeichnung schon trägt, am Ritzel auf $15$\,µm und am Rad auf $1$\,µm. Die Zahndicke ändert sich im Modell deshalb kaum. Am Ritzel liegt der Fußkreis $0{,}073$\,mm tiefer, und die lineare Kopfrücknahme nimmt ab dem Durchmesser $51{,}946$\,mm $23{,}6$\,µm zurück. Sie reicht auf der Eingriffslinie $1{,}15$\,mm vor E und erfasst damit den ganzen Doppeleingriff hinter D, wo sie noch $14$\,µm beträgt. Am Rad trägt die Profil-Winkelabweichung von $+10{,}25$\,µm Material zum Kopf hin auf, und der Bogen von $0{,}21$\,mm ersetzt die Fase der Erzeugung. Das Netz trägt drei Elementreihen auf dem Bogen und mit der Kopfkappe 87\,955 Knoten statt 83\,650. Die Normpunkte A bis E sind die der gefertigten Paarung mit der Profilüberdeckung $1{,}155$ gegen $1{,}154$ der Nenngeometrie.

\begin{figure}[!htbp]
	\centering
	\includegraphics[width=0.95\textwidth]{kap5/vergleich_referenz_w2_12nm}
	\caption{Gefertigte Paarung gegen die Nenngeometrie über dem Wälzweg, W2, quasistatische Karte, $12$\,Nm, Größtwerte über die Zahnbreite: Verdrehwinkel, größte Zugspannung der Fußrundung, größter Kontaktdruck und die Kraft des Nebenzahns.}
	\label{fig:vergleich_referenz_w2_12nm}
\end{figure}

\begin{table}[!htbp]
	\caption{Gefertigte Paarung gegen die Nenngeometrie bei $12$\,Nm, quasistatische Karte, relative Änderung in \%: größte Zugspannung $\sigma_{\mathrm{1}}$ der Fußrundung im Maximum über dem Wälzweg und an B, C und E, größter Kontaktdruck $p_{\mathrm{max}}$ im Maximum und an C, Verdrehwinkel an C.}
	\label{tab:gemessen_vergleich}
	\centering
	\footnotesize
	\setlength{\tabcolsep}{4pt}
	\begin{tabular}{L{1.2cm}L{1.7cm}L{1.5cm}L{1.5cm}L{1.5cm}L{1.9cm}L{1.6cm}L{1.9cm}}
		Stufe & $\sigma_{\mathrm{1}}$ max & $\sigma_{\mathrm{1}}$ an B & $\sigma_{\mathrm{1}}$ an C & $\sigma_{\mathrm{1}}$ an E & $p_{\mathrm{max}}$ max & $p_{\mathrm{max}}$ an C & $\varphi_{\mathrm{1}}$ an C \\
		\hline
		\arrayrulecolor{black!30}
		W1    & $+7{,}3$            & $+22{,}6$            & $+6{,}9$             & $+20{,}5$            & $-31{,}2$        & $+24{,}5$         & $+8{,}7$              \\
		\hline
		W2    & $+3{,}3$            & $+9{,}2$             & $+4{,}8$             & $+9{,}8$             & $-14{,}0$        & $+29{,}3$         & $+10{,}0$             \\
		\hline
		W3    & $+9{,}5$            & $+19{,}6$            & $+7{,}5$             & $+20{,}7$            & $-31{,}0$        & $+25{,}1$         & $+8{,}7$              \\
		\hline
		W4    & $+6{,}2$            & $+11{,}5$            & $+2{,}2$             & $+10{,}7$            & $-20{,}9$        & $+15{,}4$         & $+9{,}6$              \\
		\hline
		\arrayrulecolor{black}
	\end{tabular}
\end{table}

\tabref{tab:gemessen_vergleich} stellt die Beanspruchung der gefertigten Paarung der Nenngeometrie bei $12$\,Nm gegenüber, \figref{fig:vergleich_referenz_w2_12nm} den Verlauf über dem Wälzweg für W2. Zwei Wirkungen überlagern sich. Die Kopfrücknahme des Ritzels mindert den Lastanteil des zweiten Zahnpaars an den Enden des Wälzwegs, die Kopfrundung des Rades in kleinerem Maß. Das Hauptpaar trägt an den Doppeleingriffspunkten einen größeren Teil der Last, und die Paarung wird nachgiebiger. Lastfrei wird das zweite Paar nicht, weil die Verformung der Paarung mit $90$ bis $150$\,µm über der Rücknahme liegt. Die Profilüberdeckung der unbelasteten Geometrie zeigt davon nichts, sie zählt die Evolventen bis zu den Kopfkreisen und steigt mit dem größeren Kopfkreis des Rades sogar leicht. Der Verdrehwinkel wächst im Maximum um $6$ bis $9$\,\% und an den Normpunkten A, B, D und E um $6$ bis $23$\,\%. Die größte Zugspannung der Fußrundung über dem Wälzweg steigt bei allen Stufen und Momenten um $3$ bis $10$\,\%. Die Kopfrücknahme des Ritzels nimmt dem Kantenkontakt am Eingriffsbeginn und am Eingriffsende die Schärfe, die Kopfrundung des Rades senkt den Druck an A und D um weitere $10$ bis $17$\,\%. Der Größtwert des Kontaktdrucks liegt bei der Nenngeometrie am Kantenkontakt der Ritzelkopfkante kurz hinter E, bei der gefertigten Paarung weiter außen im verlängerten Eingriff. Die beiden Größtwerte unterscheiden sich um $8$ bis $36$\,\%. An den festen Stellen B und E, wo die Nenngeometrie ihre höchsten Drücke hat, sinkt der Druck um $16$ bis $58$\,\%. Die Kraft des Nebenzahns im verlängerten Eingriff sinkt um $1$ bis $6$\,\%.

Am Wälzpunkt C liegt der Kontaktdruck bei $8$ und $12$\,Nm dagegen um $5$ bis $29$\,\% höher. Nach den Einzelvarianten stammt diese Erhöhung vom Ritzel, nicht von der Profil-Winkelabweichung des Rades. Die Kopfrücknahme nimmt dem Nebenzahn an C fast den ganzen Lastanteil, das Hauptpaar trägt $10$\,\% mehr, und der Verdrehwinkel wächst dort bei $12$\,Nm um $11$\,\%. Bei $16$\,Nm liegt der Kontaktdruck an C unter den plastischen Stufen um $7$ bis $16$\,\% niedriger. Die Zugspannung an C ändert sich um $-2$ bis $+8$\,\%, an B und E um $+5$ bis $+33$\,\%. Dort entlastet der weichere Kantenkontakt den Nebenzahn früher, und der Hauptzahn trägt mehr Last. Die Stufen unterscheiden sich in der Höhe, nicht im Vorzeichen. Die linearen Stufen W1 und W3 reagieren am stärksten, die plastischen Stufen W2 und W4 mindern den Zuwachs der Fußspannung um ein Sechstel bis zwei Drittel.

Nach den Einzelvarianten liegt die Änderung in der Mikrogeometrie des Ritzels. Das Ritzel allein erhöht bei $12$\,Nm die Zugspannung an C unter W1 um $7{,}6$\,\% und unter W2 um $4{,}7$\,\%. An B und E steigt sie um $7$ bis $18$\,\%, der größte Kontaktdruck sinkt um $21$ bis $32$\,\%. Das Rad allein ändert die Zugspannung an C um weniger als $1$\,\%, an B und E um $2$ bis $5$\,\% und den größten Kontaktdruck um $2$ bis $4$\,\%. Die Anteile überlagern sich an der Fußrundung bis auf einen Prozentpunkt additiv.

Der Nachweis nach Stommel et al. mit W2 auf der gefertigten Paarung findet seinen kritischen Knoten an derselben Stelle wie auf der Nenngeometrie. Er liegt in der Fußrundung zwischen T3 und T4 nahe der Stirnseite bei $z = -5{,}93$\,mm. Die zyklische Auslastung steigt von $1{,}30$, $1{,}64$ und $1{,}96$ auf $1{,}38$, $1{,}70$ und $2{,}00$. Die vorhergesagte Lastwechselzahl bis zum Anriss fällt bei $8$\,Nm von $1{,}7 \cdot 10^{5}$ auf $6{,}5 \cdot 10^{4}$ und bei $12$\,Nm von $4{,}3 \cdot 10^{3}$ auf $2{,}3 \cdot 10^{3}$, bei $16$\,Nm bleibt es beim sofortigen Versagen. Der statische Nachweis besteht weiterhin an allen Fußknoten. Mit der Karte DY2, die der Dehnrate des Prüfstands nach \secref{sec:stufen_raten_koerper} am nächsten liegt, sinkt die zyklische Auslastung auf $0{,}94$, $1{,}17$ und $1{,}33$. Bei $8$\,Nm besteht der Nachweis damit für $10^{7}$ Lastwechsel, und die vorhergesagte Lastwechselzahl bis zum Anriss von $3{,}5 \cdot 10^{7}$ liegt über der gemessenen Bruchlinie für $50$\,\% mit $2{,}7 \cdot 10^{7}$. Gegen die ersten sichtbaren Anrisse der beiden Räder dieser Stufe zwischen $3{,}2$ und $7{,}7 \cdot 10^{6}$ Lastwechseln (\secref{sec:versuchsprogramm}) liegt sie um den Faktor $5$ bis $11$ zu hoch. Bei $12$ und $16$\,Nm liegt sie mit $5{,}2 \cdot 10^{5}$ und $4{,}7 \cdot 10^{4}$ um die Faktoren $11$ und $29$ unter der Bruchlinie. Die Hill-Erweiterung verhält sich wie auf der Nenngeometrie. Unter W3 versagt der statische Nachweis an den Fußknoten jenseits der Isotherme. Unter W4 liegt der kritische Knoten in der Fußrundung zwischen T2 und T3 an derselben Stelle wie zuvor, mit einer zyklischen Auslastung von $2{,}57$ gegen $2{,}45$ bei $8$\,Nm.

% Herkunft:
% Absatz 1 und 2: 80_output/fe/kst_e/vergleich_fussrundung_AtoE/analyse.md Kopf und Abschnitt 1 (Konturen);
%   ADR-116 Punkt 52 und 53 (architecture_decisions.md); Tabelle 1: analyse.md Abschnitt 1 und
%   fussrundung_sweep/bericht.md (Spalte Δ der Literaturwerte und Optima).
% Tabelle 2 und Absatz "Reihenfolge": analyse.md Abschnitt 3.1 (Δσ1 an C, ΔDrehung an C) und 3.4 (eben gegen 3D).
% Absatz "Patentkonturen": analyse.md 3.1 (Spannen Δσ1 über A bis E, Δ|σ3|, ΔCPRESS) und Abschnitt 4 (16 Nm, Drehung und Kopfweg).
% Absatz "Ellipsen": analyse.md 3.1, 3.5 und 3.6 (Ort des Größtwerts r 24,48 bis 24,60 mm gegen 24,83 mm),
%   Abschnitt 4 (|σ3| W3_QS +46/+47 %, W3_DY1 +51 %); known_limits.md FE-25 (Netzgrenze).
% Absatz "Werkstoffstufen, Dehnraten und Last": analyse.md 3.1 (W2 gegen W1, W4 gegen W3), 3.2 (Raten), 3.3 (Lasten).
% Letzter Absatz: ADR-116 Punkt 53 (Kette: Grundserie und Vollserie des bionischen Optimums).
% Bildwunsch: Konturbild der sechs Fußrundungen gegen die Trochoide im Stirnschnitt (gearcore.fillet, noch zu zeichnen, T5).
% Nomenklatur: S_L (Lückenweite als Sehne am Übergangsdurchmesser, Patent DE 10 2008 045 318 B3) fehlt noch in nomenklatur.yaml.
